% =====================================================================
%  Fiche d’activités : les 4 activités du chapitre 7 (physique, sciences mathématiques) — deux pages A4
%  L’énergie thermique : le transfert thermique (activités 1 et 2 page 1, activités 3 et 4 page 2)
%  (sans situation de départ, sans zone de nom, sans bandeau de titre)
%  À coller dans le cahier. Compilation : xelatex fiche_unique.tex (deux passes)
% =====================================================================
\documentclass[10pt,a4paper]{article}
\usepackage[a4paper,top=.75cm,bottom=.6cm,left=.8cm,right=.8cm]{geometry}
\usepackage{fontspec}
\usepackage{polyglossia}
\setmainlanguage{french}
\usepackage{amsmath}
\usepackage{unicode-math}
\setmainfont{TeX Gyre Pagella}
\setsansfont{Poppins}[Path=fonts/, Extension=.ttf, UprightFont=*-Regular, BoldFont=*-Bold, ItalicFont=*-Italic, BoldItalicFont=*-BoldItalic]
\setmathfont{Latin Modern Math}
\setmathfont[range={it/{latin,Latin,greek,Greek}}]{texgyrepagella-italic.otf}
\setmathfont[range={up/{latin,Latin,num}}]{texgyrepagella-regular.otf}
\usepackage[table]{xcolor}
\usepackage{graphicx,tikz,array}
\usepackage[most]{tcolorbox}
\graphicspath{{images/}}
\usepackage{enumitem}
\usepackage{microtype}
\usepackage[newcommands]{ragged2e}
\setlength{\RaggedRightRightskip}{0pt plus 8em}
\usepackage[version=4]{mhchem}
\usepackage{siunitx}
\sisetup{locale=FR, per-mode=power, inter-unit-product=\ensuremath{\mathord{\cdot}}, group-minimum-digits=5}
\pagestyle{empty}
\setlength{\parindent}{0pt}
\setlength{\parskip}{0pt}
\relpenalty=10000 \binoppenalty=10000
\emergencystretch=3em

\definecolor{cInk}{HTML}{1B2A41}  \definecolor{cInk2}{HTML}{51627A}
\definecolor{cBlue}{HTML}{1F5FA8} \definecolor{cOrange}{HTML}{F07F1E}
\definecolor{cTeal}{HTML}{0E8F8A} \definecolor{cSoft}{HTML}{F4F7FB}
\definecolor{cRed}{HTML}{D62839}  \definecolor{cGreen}{HTML}{1F9D74}
\color{cInk}

\newcommand{\kw}[1]{{\sffamily\bfseries\color{cBlue}#1}}
\NewDocumentCommand{\docbadge}{m}{\tikz[baseline=(b.base)]\node[fill=cBlue,text=white,rounded corners=4pt,inner xsep=3pt,inner ysep=1.4pt,font=\sffamily\bfseries\tiny](b){Doc #1};}
% document : {largeur}{n}{légende}{image}
\newcommand{\doc}[4]{%
  \begin{minipage}[t]{#1}\vspace{0pt}\centering
    \includegraphics[width=\linewidth]{#4}\par\vspace{.5mm}
    {\sffamily\scriptsize\docbadge{#2}\enspace\color{cInk2}#3}
  \end{minipage}}
\NewDocumentCommand{\pastille}{m m}{\tikz[baseline=(n.base)]\node[circle,fill=#1,text=white,inner sep=0pt,minimum size=3.9mm,font=\sffamily\bfseries\tiny](n){#2};}
\setlist[enumerate,1]{label={\pastille{cBlue}{\arabic*}}, leftmargin=6mm, labelsep=1.8mm, itemsep=2pt, topsep=2pt, parsep=0pt}
% séparateur discret entre les activités
\newcommand{\sep}[1]{\par\vspace{2.5mm}%
  \noindent\tikz[baseline=(t.base)]\node[fill=cTeal,text=white,rounded corners=3pt,inner xsep=4pt,inner ysep=1.5pt,font=\sffamily\bfseries\scriptsize](t){Activité #1};%
  \hspace{2mm}{\color{cTeal!70}\leaders\hrule height .7pt\hfill\kern0pt}\par\vspace{1mm}}

% ---------- tableau de mesures compact ----------
\newcommand{\mesures}[3][0.8cm]{{\renewcommand{\arraystretch}{1.2}\setlength{\arrayrulewidth}{.5pt}\arrayrulecolor{cTeal!55}\setlength{\tabcolsep}{3pt}%
  \begin{tabular}{|>{\sffamily\scriptsize\bfseries\color{cTeal}}l|*{#2}{>{\centering\arraybackslash}p{#1}|}}\hline #3\end{tabular}}}

\newcommand{\vect}[1]{\overrightarrow{#1}}
\newcommand{\Ec}{E_{\mathrm{C}}}
\newcommand{\dC}[1]{\qty{#1}{\degreeCelsius}}
\newcommand{\cu}{\unit{J.kg^{-1}.\degreeCelsius^{-1}}}
\newcommand{\tfus}{\theta_{\mathrm{fus}}}
\newcommand{\info}[1]{\par\vspace{.5mm}{\setlength{\fboxsep}{2pt}\colorbox{cSoft}{\parbox{\dimexpr\linewidth-4pt}{\RaggedRight #1}}}\par\vspace{.5mm}}
\begin{document}
\small\RaggedRight
\setlength{\parskip}{0pt}
% =====================================================================
\sep{1}
\begin{minipage}[t]{0.64\linewidth}\vspace{0pt}\RaggedRight
Des thermoplongeurs identiques sont reliés à des \textbf{joulemètres} : l’énergie mesurée est entièrement transférée au corps sous forme de chaleur. Corps pris à $\theta_i = \dC{20,0}$.\par\vspace{.5mm}
\textbf{Doc 1.} Eau chauffée jusqu’à \dC{40,0}.\enspace
\mesures[0.75cm]{2}{%
Masse (\unit{g}) & 150 & 300\\ \hline
Énergie (\unit{kJ}) & 12,6 & 25,1\\ \hline}\par\vspace{.5mm}
\textbf{Doc 2.} \qty{200}{g} d’eau.\enspace
\mesures[0.75cm]{3}{%
$\Delta\theta$ (\unit{\degreeCelsius}) & 10,0 & 20,0 & 30,0\\ \hline
Énergie (\unit{kJ}) & 8,37 & 16,7 & 25,1\\ \hline}\par\vspace{.5mm}
\textbf{Doc 3.} Même énergie (\qty{5,02}{kJ}) pour \qty{200}{g} d’eau et \qty{200}{g} de fer.\enspace
\mesures[0.7cm]{2}{%
Corps & eau & fer\\ \hline
$\Delta\theta$ (\unit{\degreeCelsius}) & 6,0 & 55,8\\ \hline}
\info{\textbf{Information.} $Q$ : chaleur \emph{reçue} par un corps ($Q > 0$ s’il reçoit de l’énergie, $Q < 0$ s’il en cède). Sans changement d’état, $Q$ est égale au produit de la masse $m$, de $\Delta\theta = \theta_f - \theta_i$ et d’un coefficient $c$ qui ne dépend que de la nature du corps : sa \textbf{chaleur massique}, en \cu{} (eau : 4\,185 ; fer : 450). $C = m\cdot c$ : \textbf{capacité thermique}.}
\end{minipage}\hfill
\begin{minipage}[t]{0.34\linewidth}\vspace{0pt}\centering
\doc{\linewidth}{1}{Thermoplongeurs et joulemètres}{doc1_thermoplongeurs}\par\vspace{1mm}
\doc{\linewidth}{3}{Eau et fer : même énergie}{doc3_eau_fer}
\end{minipage}
\begin{enumerate}
  \item \kw{Établir.} Docs 1 et 2 : montrer que $Q$ est proportionnelle à la masse $m$, puis à $\Delta\theta$.
  \item \kw{Interpréter.} Doc 3 : combien de fois $\Delta\theta$ du fer est-il plus grand que celui de l’eau ? Comparer avec le rapport des chaleurs massiques.
  \item \kw{Écrire.} À l’aide de l’information, écrire $Q$. Que représente $c$ si $m = \qty{1}{kg}$ et $\Delta\theta = \dC{1}$ ? Que représente $C = m\cdot c$ ? Unités ?
  \item \kw{Relier.} Signe de $Q$ quand le corps se refroidit ? Pour un corps au repos qui n’échange pas de travail, que vaut $\Delta U$ (premier principe) ?
\end{enumerate}

% =====================================================================
\sep{2}
\begin{minipage}[t]{0.72\linewidth}\vspace{0pt}\RaggedRight
\textbf{Doc 4. Le calorimètre.} Enceinte thermiquement isolée : les corps qu’il contient échangent de la chaleur jusqu’à l’\textbf{équilibre thermique}, où ils ont tous la même température, constante. L’énergie cédée par les corps qui se refroidissent est entièrement reçue par ceux qui s’échauffent ; le bilan ne dépend que des états initial et final. Le calorimètre et ses accessoires reçoivent $Q = \mu\cdot(\theta_f - \theta_i)$, où $\mu$ est leur \textbf{capacité thermique}.\par\vspace{.4mm}
\textbf{Étape 1.} $m_1 = \qty{250}{g}$ d’eau à $\theta_1 = \dC{16,5}$ ; on verse $m_2 = \qty{200}{g}$ d’eau à $\theta_2 = \dC{52,0}$ ; équilibre : $\theta_{e1} = \dC{31,5}$.\par\vspace{.4mm}
\textbf{Étape 2.} Le calorimètre vidé reçoit $m_e = \qty{300}{g}$ d’eau à $\theta_0 = \dC{19,0}$ ; on y plonge $m_{\mathrm{Cu}} = \qty{250}{g}$ de cuivre à $\theta_m = \dC{96,0}$ ; équilibre : $\theta_{e2} = \dC{24,2}$. $c_e = \qty{4185}{J.kg^{-1}.\degreeCelsius^{-1}}$.
\begin{enumerate}
  \item \kw{Définir.} D’après le doc 4, quand des corps sont-ils en équilibre thermique ? Que vaut la somme $Q_1 + Q_2 + \dots$ des chaleurs échangées (\textbf{équation calorimétrique}) ?
  \item \kw{Calculer.} Étape 1 : exprimer les chaleurs échangées ; en déduire $\mu$.
  \item \kw{Exprimer.} Étape 2 : exprimer les chaleurs échangées par \{calorimètre + eau\} et par le cuivre.
  \item \kw{Déterminer.} En déduire $c_{\mathrm{Cu}}$ (avec $\mu$ de la question 2) ; comparer à \qty{385}{J.kg^{-1}.\degreeCelsius^{-1}}.
\end{enumerate}
\end{minipage}\hfill
\doc{0.26\linewidth}{4}{Calorimètre (coupe)}{doc4_calorimetre}

\newpage
% =====================================================================
\sep{3}
\begin{minipage}[t]{0.64\linewidth}\vspace{0pt}\RaggedRight
\textbf{Doc 5.} On chauffe à puissance constante \qty{100}{g} de glace prise à \dC{-10} et on relève sa température.\par\vspace{.4mm}
\textbf{Chaleur latente.} À pression constante, un corps pur change d’état à température constante ($\tfus = \dC{0}$ pour la glace). La chaleur échangée est proportionnelle à la masse $m$ qui change d’état : $Q = m\cdot L$, $L$ : \textbf{chaleur latente}. La fusion se fait en recevant de la chaleur ; la solidification, inverse, cède exactement la même quantité.\par\vspace{.4mm}
\textbf{Doc 6.} Calorimètre : $\mu = \qty{75}{J.\degreeCelsius^{-1}}$, $m_1 = \qty{350}{g}$ d’eau à $\theta_1 = \dC{24,0}$. On y introduit $m_2 = \qty{50}{g}$ de glace fondante (\dC{0,0}). Après fusion complète : $\theta_e = \dC{11,4}$. $c_e = \qty{4185}{J.kg^{-1}.\degreeCelsius^{-1}}$.
\end{minipage}\hfill
\begin{minipage}[t]{0.34\linewidth}\vspace{0pt}\centering
\doc{\linewidth}{5}{Glace chauffée}{doc5_courbe}\par\vspace{1mm}
\doc{0.62\linewidth}{6}{Glace dans le calorimètre}{doc6_glacon}
\end{minipage}
\begin{enumerate}
  \item \kw{Exploiter.} Doc 5 : décrire les trois parties de la courbe. La glace reçoit-elle de l’énergie pendant le palier ?
  \item \kw{Définir.} D’après $Q = m\cdot L$, que représente $L$ si $m = \qty{1}{kg}$ ? Unité ? Relation entre $L_s$ (solidification) et $L_f$ (fusion) ?
  \item \kw{Décrire.} Doc 6 : décrire les deux étapes subies par la glace ; exprimer la chaleur totale $Q_g$ qu’elle reçoit.
  \item \kw{Exprimer.} Exprimer la chaleur $Q_{ce}$ cédée par le calorimètre et l’eau.
  \item \kw{Calculer.} Appliquer l’équation calorimétrique pour calculer $L_f$.
\end{enumerate}

% =====================================================================
\sep{4}
\begin{minipage}[t]{0.64\linewidth}\vspace{0pt}\RaggedRight
\textbf{Doc 7.} Un capteur à surface noircie (thermomètre à réservoir noirci) est placé sous une cloche où l’on a fait le vide. Une lampe extérieure (visible et infrarouge) l’éclaire : dès l’allumage, sa température s’élève.
\info{\textbf{Information.} La lumière visible et l’infrarouge sont des \textbf{ondes électromagnétiques} : elles transportent de l’énergie, se propagent en ligne droite et même dans le vide. Ce mode de transfert s’appelle le \textbf{rayonnement}. Un corps qui absorbe un rayonnement voit son énergie interne augmenter.}
\textbf{Doc 8.} \textbf{Conduction} : dans un solide, de proche en proche, sans déplacement de matière. \textbf{Convection} : dans un liquide ou un gaz, la matière chaude se déplace. Ces deux modes ont besoin de matière. Une surface sombre et mate \textbf{absorbe} la plus grande partie du rayonnement reçu ; une surface brillante ou blanche le \textbf{réfléchit}. Une \textbf{vitre} laisse passer la lumière visible mais arrête une grande partie de l’infrarouge émis par les objets chauds.
\end{minipage}\hfill
\doc{0.34\linewidth}{7}{Capteur sous une cloche à vide}{doc7_cloche}
\begin{enumerate}
  \item \kw{Interpréter.} Doc 7 : pourquoi l’énergie ne peut-elle pas arriver au capteur par l’air ? Nommer le mode de transfert.
  \item \kw{Comparer.} Doc 8 : quelle différence essentielle y a-t-il entre le rayonnement et la conduction ou la convection ?
  \item \kw{Expliquer.} Pourquoi les tubes d’un capteur solaire thermique sont-ils peints en noir mat et placés sous une vitre ?
\end{enumerate}
\end{document}
